\begin{frame}{지도학습 내부: 판별적 vs 생성적}
  지도학습 안에서도 다시 두 갈래가 갈린다:
  \begin{columns}[T]
    \begin{column}{0.5\textwidth}
      \kb{판별적 (discriminative)}
      \begin{itemize}
        \item $P(y \mid x)$ --- ``$x$ 가 주어졌을 때 $y$ 의 확률''을
          \textbf{곧바로} 학습.
        \item 예: 회귀 · 로지스틱회귀 · SVM · 신경망.
      \end{itemize}
    \end{column}
    \begin{column}{0.5\textwidth}
      \kb{생성적 (generative)}
      \begin{itemize}
        \item $P(x \mid y)$ --- ``각 클래스가 데이터를 \textbf{어떻게
          만들어내는지}''를 먼저 학습.
        \item 그 다음 \textbf{베이즈 정리로 뒤집어} $P(y \mid x)$ 를
          얻는다. 예: Week 3의 나이브베이즈·GDA.
      \end{itemize}
    \end{column}
  \end{columns}
  \vskip0.5em
  같은 문제를 \textbf{정반대 방향}에서 접근하는 두 철학의 차이는
  Week 3에서 자세히 다룬다.
  \vskip0.3em
  {\footnotesize\color{ksagray} 여기서 ``생성적''은 \textbf{분류}를
  위해 $P(x\mid y)$를 학습하는 것 --- 이미지·텍스트를 새로 만들어내는
  Week 15의 \textbf{생성형 모델}(VAE$\cdot$GAN$\cdot$Diffusion, LLM
  포함)과는 별개의 과제이며, 지도학습 분류에 속하지 않는다. 다만
  ``데이터가 어떻게 만들어지는지''를 모델링한다는 철학은 같아서,
  이 관점이 그 생성형 모델들로 이어지는 첫 걸음이다.}
\end{frame}
